\documentclass[12pt]{book}
\usepackage[paperwidth=17cm, paperheight=24cm, margin=0pt]{geometry}
\usepackage{xcolor}
\usepackage{tikz}
\usepackage{fontenc}
\usepackage[utf8]{inputenc}
\usepackage{amsmath}
\usetikzlibrary{calc, patterns, decorations.pathmorphing}

% ── Paleta de colores ─────────────────────────────────────────────
\definecolor{ugazul}    {RGB}{0,  51, 102}   % azul profundo UG
\definecolor{ugclaro}   {RGB}{0,  90, 160}   % azul medio
\definecolor{ugacento}  {RGB}{200,160,  0}   % dorado/ocre
\definecolor{ugfondo}   {RGB}{245,245,240}   % crema muy suave
\definecolor{uglinea}   {RGB}{180,195,215}   % azul grisáceo para grid
\definecolor{ugtexto}   {RGB}{20,  30,  50}  % casi negro azulado

\pagestyle{empty}

\begin{document}
\begin{tikzpicture}[remember picture, overlay]

  %% ══════════════════════════════════════════════════════════════
  %% 1. FONDO BASE – crema
  %% ══════════════════════════════════════════════════════════════
  \fill[ugfondo] (current page.south west) rectangle (current page.north east);

  %% ══════════════════════════════════════════════════════════════
  %% 2. PANEL LATERAL IZQUIERDO – franja azul oscura
  %% ══════════════════════════════════════════════════════════════
  \fill[ugazul]
    (current page.south west)
    rectangle
    ([xshift=3.8cm]current page.north west);

  %% ══════════════════════════════════════════════════════════════
  %% 3. BANDA DORADA – línea de acento entre franja y contenido
  %% ══════════════════════════════════════════════════════════════
  \fill[ugacento]
    ([xshift=3.8cm]current page.south west)
    rectangle
    ([xshift=4.15cm]current page.north west);

  %% ══════════════════════════════════════════════════════════════
  %% 4. GRID MATEMÁTICO – papel cuadriculado sutil en zona crema
  %% ══════════════════════════════════════════════════════════════
  \begin{scope}
    \clip ([xshift=4.15cm]current page.south west)
          rectangle (current page.north east);
    % líneas verticales
    \foreach \x in {4.5,5.0,...,17}{
      \draw[uglinea, line width=0.3pt, opacity=0.5]
        ([xshift=\x cm]current page.south west) --
        ([xshift=\x cm]current page.north west);
    }
    % líneas horizontales
    \foreach \y in {0,0.5,...,24}{
      \draw[uglinea, line width=0.3pt, opacity=0.5]
        ([yshift=\y cm]current page.south west) --
        ([yshift=\y cm]current page.south east);
    }
  \end{scope}

  %% ══════════════════════════════════════════════════════════════
  %% 5. BLOQUE SUPERIOR AZUL – cabecera con título
  %% ══════════════════════════════════════════════════════════════
  \fill[ugazul]
    ([xshift=4.15cm, yshift=-7.5cm]current page.north west)
    rectangle
    (current page.north east);

  %% ══════════════════════════════════════════════════════════════
  %% 6. BANDA DORADA INFERIOR del bloque azul superior
  %% ══════════════════════════════════════════════════════════════
  \fill[ugacento]
    ([xshift=4.15cm, yshift=-7.5cm]current page.north west)
    rectangle
    ([yshift=-7.85cm]current page.north east);

  %% ══════════════════════════════════════════════════════════════
  %% 9. TEXTO DEL TÍTULO (en bloque azul superior)
  %% ══════════════════════════════════════════════════════════════
  \node[anchor=north west, text=white, align=left, inner sep=0pt,
        text width=11cm]
    at ([xshift=5.0cm, yshift=-1.0cm]current page.north west)
    {%
      \fontsize{9}{11}\selectfont\textit{\textcolor{ugacento}{\'Algebra Lineal}}\\[6pt]
      \fontsize{26}{30}\selectfont\bfseries
      \'Algebra\\[2pt] Lineal\\[8pt]
      \fontsize{10}{13}\selectfont\normalfont\textcolor{uglinea}
      {Notas del curso}%
    };

  %% ══════════════════════════════════════════════════════════════
  %% 10. FÓRMULA DECORATIVA en bloque azul
  %% ══════════════════════════════════════════════════════════════
  \node[anchor=south east, text=white, opacity=0.18,
        font=\fontsize{42}{42}\selectfont, inner sep=0pt]
    at ([xshift=-0.6cm, yshift=-0.3cm]current page.north east)
    {$A\mathbf{x}=\lambda\mathbf{x}$};

  %% ══════════════════════════════════════════════════════════════
  %% 11. TEXTO EN FRANJA LATERAL – rotado 90°
  %% ══════════════════════════════════════════════════════════════
  \node[rotate=90, anchor=south, text=white, opacity=0.55,
        font=\fontsize{7}{9}\selectfont\scshape, inner sep=0pt]
    at ([xshift=1.9cm, yshift=3.45cm]current page.south west)
    {Universidad de Guanajuato \quad$\cdot$\quad Departamento de Matem\'aticas};

  %% ══════════════════════════════════════════════════════════════
  %% 12. ZONA CENTRAL – contenido principal
  %% ══════════════════════════════════════════════════════════════

  % Rectángulo decorativo de énfasis para autores
  \fill[ugazul, opacity=0.07]
    ([xshift=4.7cm, yshift=-15.5cm]current page.north west)
    rectangle
    ([xshift=-1cm, yshift=-20.5cm]current page.north east);

  \draw[ugazul, opacity=0.25, line width=0.6pt]
    ([xshift=4.7cm, yshift=-15.5cm]current page.north west)
    rectangle
    ([xshift=-1cm, yshift=-20.5cm]current page.north east);

  % Línea de acento dorada izquierda del bloque autores
  \fill[ugacento]
    ([xshift=4.7cm, yshift=-15.5cm]current page.north west)
    rectangle
    ([xshift=5.05cm, yshift=-20.5cm]current page.north west);

  % Etiqueta "Autores"
  \node[anchor=north west, text=ugacento,
        font=\fontsize{7.5}{9}\selectfont\scshape, inner sep=0pt]
    at ([xshift=5.5cm, yshift=-15.9cm]current page.north west)
    {Autores};

  % Nombres autores
  \node[anchor=north west, text=ugtexto, align=left, inner sep=0pt]
    at ([xshift=5.5cm, yshift=-16.7cm]current page.north west)
    {%
      \fontsize{13}{17}\selectfont\bfseries
      Carlos Alberto Rom\'an Rodr\'iguez\\[3pt]
      Ricardo Le\'on Mart\'inez%
    };

  % Etiqueta "Profesora"
  \node[anchor=north west, text=ugacento,
        font=\fontsize{7.5}{9}\selectfont\scshape, inner sep=0pt]
    at ([xshift=5.5cm, yshift=-18.8cm]current page.north west)
    {Profesora del curso};

  % Nombre profesora
  \node[anchor=north west, text=ugtexto, align=left, inner sep=0pt]
    at ([xshift=5.5cm, yshift=-19.5cm]current page.north west)
    {%
      \fontsize{12}{16}\selectfont
      Claudia Reynoso Alc\'antara%
    };

  %% ══════════════════════════════════════════════════════════════
  %% 13. BLOQUE INFERIOR – universidad e institución
  %% ══════════════════════════════════════════════════════════════
  \fill[ugazul]
    (current page.south west)
    rectangle
    ([yshift=3.2cm]current page.south east);

  \fill[ugacento]
    ([yshift=3.2cm]current page.south west)
    rectangle
    ([yshift=3.55cm]current page.south east);

  % Universidad
  \node[anchor=west, text=white, inner sep=0pt,
        font=\fontsize{11}{13}\selectfont\bfseries]
    at ([xshift=4.5cm, yshift=2.3cm]current page.south west)
    {Universidad de Guanajuato};

  % Departamento
  \node[anchor=west, text=uglinea, inner sep=0pt,
        font=\fontsize{8.5}{11}\selectfont]
    at ([xshift=4.5cm, yshift=1.4cm]current page.south west)
    {Departamento de Matem\'aticas};

  % Año
  \node[anchor=east, text=ugacento, inner sep=0pt,
        font=\fontsize{10}{12}\selectfont\bfseries]
    at ([xshift=-1.2cm, yshift=1.85cm]current page.south east)
    {2026};

  %% ══════════════════════════════════════════════════════════════
  %% 14. MATRIZ DECORATIVA – fondo zona media derecha
  %% ══════════════════════════════════════════════════════════════
  \node[anchor=south east, text=ugazul, opacity=0.09,
        font=\fontsize{52}{52}\selectfont, inner sep=0pt]
    at ([xshift=-1cm, yshift=3.8cm]current page.south east)
    {$\begin{pmatrix} a & b \\ c & d \end{pmatrix}$};

\end{tikzpicture}
\end{document}
